\documentclass[11pt]{article}
\usepackage[margin=1in]{geometry}
\usepackage{amsmath,amssymb,amsthm,mathtools}
\usepackage{booktabs}
\usepackage{microtype}
\usepackage[colorlinks=true,linkcolor=blue,urlcolor=blue,citecolor=blue]{hyperref}

\newtheorem{theorem}{Theorem}
\newtheorem{proposition}[theorem]{Proposition}
\newtheorem{corollary}[theorem]{Corollary}
\newtheorem{lemma}[theorem]{Lemma}
\newcommand{\R}{\mathbb R}
\newcommand{\Sn}{S_n}
\newcommand{\dist}{\operatorname{dist}}
\newcommand{\down}{\mathord{\downarrow}}

\title{Phase retrieval is strictly weaker than universal sorting once there are three rows}
\author{Candidate full implication-level answer to the open question in arXiv:2306.13111}
\date{11 August 2026}

\begin{document}
\maketitle

\noindent\fbox{\parbox{0.94\linewidth}{\textbf{Status: candidate full result,
likely valid.}  For every number of rows \(n\ge2\), an injective sorting
encoder \(\beta_A\) forces the phase-retrieval map \(\alpha_A\) to be
injective, and the same implication holds with unchanged admissible
bi-Lipschitz constants.  The converse is an equivalence for \(n=2\), as proved
in the source, but is strict for every \(n\ge3\) and \(d\ge2\).  More sharply,
every minimally redundant real phase-retrieval frame, with \(D=2d-1\)
measurements, fails to be a universal sorting key for every \(n\ge3\).
Consequently a universal key for \(n\ge3\) requires \(D\ge2d\).}}

\section{\texorpdfstring{\hspace{0.35em}}{}The source question}

Let \(A=[a_1|\cdots|a_D]\in\R^{d\times D}\).  The source compares
\[
 \alpha_A(x)=\bigl(|\langle x,a_j\rangle|\bigr)_{j=1}^D
\]
on \(\R^d/\{\pm1\}\) with the columnwise sorting encoder
\[
 \beta_A(X)=\down(XA),\qquad X\in\R^{n\times d},
\]
on \(\R^{n\times d}/\Sn\).  The matrix \(A\) is called a
\emph{universal key} when \(\beta_A\) separates row-permutation orbits.
Theorem~1 of the source proves that, for \(n=2\), \(A\) is a universal key if
and only if it is a real phase-retrieval frame.  The conclusion asks for the
relationship between \(\alpha_A\) and \(\beta_A\) when \(n>2\)
\cite[Conclusion]{BT}.

The result below completely determines the implication relation between the
two injectivity properties.  It also gives a quantitative one-way relation
and a new template-count obstruction.

\section{\texorpdfstring{\hspace{0.35em}}{}Idea of the proof}

There are two simple but complementary ideas.

First, insert a signal \(x\) into the \(n\)-row point cloud
\[
 X_x=(x,-x,0,\ldots,0)^T.
\]
Its permutation-orbit distance is exactly \(\sqrt2\) times the sign-orbit
distance of \(x\), while sorting each projected column merely places
\(-|\langle x,a_j\rangle|\) and \(+|\langle x,a_j\rangle|\) around zeros.
Thus phase retrieval sits as a scaled isometric slice of the sorting problem.

Second, let \(A\) be a full-spark frame with the minimal \(2d-1\) columns.
Partition the columns into blocks of sizes \(d-1,d-1,1\).  Normals to the
first two blocks can be scaled so that their sum is orthogonal to the last
column.  This produces three nonzero vectors \(r_1,r_2,r_3\) with sum zero
such that every measurement vector is orthogonal to at least one \(r_i\).
Along such a measurement, the three projections of
\(\{r_1,r_2,r_3\}\) have the form \(\{0,t,-t\}\), so they are unchanged after
negating the entire point cloud.  The two clouds are nevertheless not row
permutations.  This is a three-point switching component tailored to every
minimal phase-retrieval frame.

\section{\texorpdfstring{\hspace{0.35em}}{}The phase-retrieval slice}

Write
\[
 d_{\pm}(x,y)=\min\{\|x-y\|,\|x+y\|\}
\]
and
\[
 d_{\Sn}(X,Y)=\min_{P\in\Sn}\|X-PY\|_F.
\]
For \(n\ge2\), let \(X_x\in\R^{n\times d}\) have first two rows \(x,-x\)
and all remaining rows zero.

\begin{lemma}\label{lem:slice}
For all \(x,y\in\R^d\),
\[
 d_{\Sn}(X_x,X_y)=\sqrt2\,d_{\pm}(x,y)
\]
and
\[
 \|\beta_A(X_x)-\beta_A(X_y)\|_F
 =\sqrt2\,\|\alpha_A(x)-\alpha_A(y)\|_2.
\]
\end{lemma}

\begin{proof}
The squared matching cost between the row multisets
\(\{x,-x,0,\ldots,0\}\) and \(\{y,-y,0,\ldots,0\}\) equals
\[
 2\|x\|^2+2\|y\|^2
 -2\max_{P\in\Sn}\sum_{i=1}^n\langle (X_x)_i,(PX_y)_i\rangle.
\]
The maximum inner-product sum is \(2|\langle x,y\rangle|\), achieved by
matching the two nonzero antipodal pairs with the favorable sign.  This gives
\[
 d_{\Sn}(X_x,X_y)^2
 =2\bigl(\|x\|^2+\|y\|^2-2|\langle x,y\rangle|\bigr)
 =2d_{\pm}(x,y)^2.
\]

For the \(j\)-th column, putting \(s=\langle x,a_j\rangle\) gives
\[
 \down(X_xa_j)=(-|s|,0,\ldots,0,|s|)^T
\]
when sorting is increasing; decreasing order only reverses the vector.
The squared difference of the \(j\)-th sorted columns for \(x\) and \(y\) is
twice the squared difference of the corresponding magnitudes.  Summing over
\(j\) proves the second identity.
\end{proof}

\begin{proposition}\label{prop:oneway}
Let \(n\ge2\).  If \(\beta_A\) is injective on
\(\R^{n\times d}/\Sn\), then \(\alpha_A\) is injective on
\(\R^d/\{\pm1\}\).

More quantitatively, if
\[
 c\,d_{\Sn}(X,Y)\le
 \|\beta_A(X)-\beta_A(Y)\|_F
 \le C\,d_{\Sn}(X,Y)
\]
for all \(X,Y\), then
\[
 c\,d_{\pm}(x,y)\le
 \|\alpha_A(x)-\alpha_A(y)\|_2
 \le C\,d_{\pm}(x,y)
\]
for all \(x,y\).  Hence the optimal lower constant of \(\alpha_A\) is at
least that of \(\beta_A\), and its optimal upper constant is at most that of
\(\beta_A\).
\end{proposition}

\begin{proof}
Restrict \(\beta_A\) to the point clouds \(X_x\) and apply
Lemma~\ref{lem:slice}.  If \(\alpha_A(x)=\alpha_A(y)\), then
\(\beta_A(X_x)=\beta_A(X_y)\); injectivity makes the row multisets equal,
which forces \(y=\pm x\).
\end{proof}

\section{Strictness for every \texorpdfstring{\(n\ge3\)}{n at least 3}}

We use the standard real complement-property characterization: a real frame
does phase retrieval if and only if every partition of its columns has a
spanning side \cite{BCE}.  With exactly \(2d-1\) columns, this is equivalent
to full spark.  Indeed, full spark makes the side with at least \(d\) columns
span; conversely, a dependent \(d\)-set and its \((d-1)\)-element complement
violate the complement property.

\begin{theorem}\label{thm:strict}
Let \(d\ge2\), \(n\ge3\), and
\(A=[a_1|\cdots|a_{2d-1}]\in\R^{d\times(2d-1)}\) be full spark.  Then
\(\beta_A\) is not injective on \(\R^{n\times d}/\Sn\).

Consequently every minimally redundant real phase-retrieval frame fails to
be a universal key as soon as \(n\ge3\).
\end{theorem}

\begin{proof}
Partition the column indices into disjoint sets
\[
 I\mathbin{\dot\cup}J\mathbin{\dot\cup}\{k\}
 =[2d-1],\qquad |I|=|J|=d-1.
\]
Full spark implies that the columns in each of \(I\) and \(J\) span a
hyperplane.  Choose nonzero normals
\[
 u\perp\operatorname{span}\{a_i:i\in I\},\qquad
 v\perp\operatorname{span}\{a_j:j\in J\}.
\]
Both \(\langle a_k,u\rangle\) and \(\langle a_k,v\rangle\) are nonzero;
otherwise \(a_k\) together with one of the \((d-1)\)-blocks would be a
dependent \(d\)-set.  Define
\[
 r_1=\langle a_k,v\rangle u,\qquad
 r_2=-\langle a_k,u\rangle v,\qquad
 r_3=-(r_1+r_2).
\]
Then
\[
 r_1+r_2+r_3=0,\qquad
 a_i\perp r_1\ (i\in I),\qquad
 a_j\perp r_2\ (j\in J),\qquad
 a_k\perp r_3.
\]
The vectors \(r_1,r_2\) are nonzero.  They are not collinear: otherwise the
two block hyperplanes would coincide and would contain the \(2d-2\) columns
indexed by \(I\cup J\), contradicting full spark.  Hence \(r_3\ne0\) as well.

Let \(X,Y\in\R^{n\times d}\) have row multisets
\[
 \{r_1,r_2,r_3,0,\ldots,0\},\qquad
 \{-r_1,-r_2,-r_3,0,\ldots,0\},
\]
respectively.  Fix a measurement column \(a_\ell\).  It is orthogonal to at
least one \(r_q\).  Since the three scalar projections sum to zero, their
multiset is
\[
 \{\langle r_1,a_\ell\rangle,
   \langle r_2,a_\ell\rangle,
   \langle r_3,a_\ell\rangle\}
 =\{0,t,-t\}
\]
for some \(t\).  Negating the three vectors leaves this multiset unchanged.
The additional zero rows agree, so every sorted column agrees and
\(\beta_A(X)=\beta_A(Y)\).

Yet \(X\) and \(Y\) are not row permutations.  If the three-element
multiset \(\{r_1,r_2,r_3\}\) equaled its negative, then an odd-cardinality
sign-invariant multiset would have to contain zero: all nonzero elements pair
as \(r,-r\) with equal multiplicities.  None of the \(r_i\) is zero, a
contradiction.
\end{proof}

\begin{corollary}\label{cor:count}
Let \(d\ge2\) and \(n\ge3\).  Every universal key
\(A\in\R^{d\times D}\) has
\[
 D\ge2d.
\]
\end{corollary}

\begin{proof}
Proposition~\ref{prop:oneway} makes \(A\) a phase-retrieval frame, so the
real complement property gives \(D\ge2d-1\).  Equality would make \(A\)
full spark and contradict Theorem~\ref{thm:strict}.
\end{proof}

\begin{corollary}\label{cor:classification}
At the level of injectivity, the relationship is:
\[
\begin{array}{c|c}
\text{regime}&\text{relationship}\\ \hline
n=2&\beta_A\text{ injective}\Longleftrightarrow
      \alpha_A\text{ injective}\quad\text{(source theorem)},\\
d=1&\beta_A\text{ injective}\Longleftrightarrow
      \alpha_A\text{ injective}\quad\text{for every }n,\\
n\ge3,\ d\ge2&\beta_A\text{ injective}\Longrightarrow
      \alpha_A\text{ injective, and the converse is false.}
\end{array}
\]
\end{corollary}

\begin{proof}
Only the one-dimensional statement remains.  A frame in \(\R\) does phase
retrieval exactly when some \(a_j\ne0\).  Sorting the corresponding nonzero
scalar multiple of the rows recovers their unordered multiset, so the same
condition makes \(\beta_A\) injective for every \(n\).

For \(n\ge3,d\ge2\), full-spark frames with \(2d-1\) columns exist, for
example Vandermonde frames.  They phase retrieve but fail universality by
Theorem~\ref{thm:strict}.
\end{proof}

\section{\texorpdfstring{\hspace{0.35em}}{}Current bounds and a sharp small case}

The 2026 follow-up by Dym, Wellershoff, Tsoukanis, Levy, and Balan proves that
every full-spark \(A\) is universal when
\[
 D\ge n(d-1)+1
\]
and gives a different logarithmic lower bound \cite[Theorems 4 and 5]{DWTLB}.
Its numerical section reports failure to find separating keys in the cases
\[
 (n,d,D)=(3,2,3),(3,3,5),(3,4,7),
\]
exactly the pattern \(n=3,D=2d-1\), but states this only as numerical
evidence \cite[Sec.~IV]{DWTLB}.  Theorem~\ref{thm:strict} proves that failure
for every \(d\ge2\), for every full-spark choice of the columns, and indeed
for every \(n\ge3\).

Combining Corollary~\ref{cor:count} with the follow-up's full-spark upper
bound gives the exact small case
\[
 D_{\min}(n=3,d=2)=4.
\]
More generally, the present result improves the universal lower bound to
\(D\ge2d\) throughout the regime \(n\ge3,d\ge2\); depending on \(n,d\), the
follow-up's logarithmic obstruction may be stronger.

\section{Proof audit, scope, and novelty}

The argument uses no genericity and no numerical inference.  The only
external structural input is the standard real complement-property theorem;
the equivalence between minimal phase retrieval and full spark is reproved
above.  The construction is homogeneous and works with repeated zero rows,
so the passage from three rows to arbitrary \(n\ge3\) is exact.

This is a full answer to the source's implication-level injectivity question,
plus a quantitative one-way result.  It does not solve the separate neural
network implementation problem mentioned in the same conclusion, nor does it
determine the optimal number of templates for every \(n,d\).

Searches through 11 August 2026 covered the source title and exact concluding
sentence, ``phase retrieval'', ``universal key'', ``sorting-based
permutation-invariant embedding'', and the later literature.  The closest
paper is arXiv:2510.22186v2, which proves broader upper and lower bounds and
records the \(n=3,D=2d-1\) failures only numerically.  The three-point
switching idea is related to classical discrete tomography
\cite{MPS}, especially in the planar case, but no source located states the
frame-uniform theorem above, the \(D\ge2d\) corollary, or the exact scaled
phase-retrieval slice.  Novelty confidence is therefore moderate.

\begin{thebibliography}{9}

\bibitem{BT}
R.~Balan and E.~Tsoukanis,
\emph{Relationships between the Phase Retrieval Problem and Permutation
Invariant Embeddings},
arXiv:2306.13111v2 (2023),
\url{https://arxiv.org/abs/2306.13111}.

\bibitem{BCE}
R.~Balan, P.~G. Casazza, and D.~Edidin,
\emph{On signal reconstruction without phase},
Applied and Computational Harmonic Analysis 20 (2006), 345--356.

\bibitem{DWTLB}
N.~Dym, M.~Wellershoff, E.~Tsoukanis, D.~Levy, and R.~Balan,
\emph{Quantitative Bounds for Sorting-Based Permutation-Invariant Embeddings},
arXiv:2510.22186v2 (2026),
\url{https://arxiv.org/abs/2510.22186}.

\bibitem{MPS}
J.~Matou\v{s}ek, A.~P\v{r}\'{\i}v\v{e}tiv\'y, and P.~\v{S}kovro\v{n},
\emph{How many points can be reconstructed from \(k\) projections?},
SIAM Journal on Discrete Mathematics 22 (2008), 1605--1623.

\end{thebibliography}

\end{document}
